\documentclass{situationist-preprint}

\title{Obligations Incurred in Childhood}
\msid{PID-07-b}
\msfamily{PID-07}
\msdiscipline{Personal Identity and the Self}
\mstype{Companion manuscript}
\mskeywords{childhood; promises; responsibility; psychological connectedness; reliance; open future}

\begin{document}
\maketitle

\begin{abstract}
PID-07-a holds that you are the same person as your childhood self, that your connectedness with it is weak, and that the obligations grounded in the relation between you and it are correspondingly weak. We examine three kinds of case in detail: promises made as a child, wrongs done as a child, and what the adult may owe to the hopes of the child she was. In each we separate the component of the obligation grounded in the relation between the adult and her childhood self, which is weak, from components with other grounds, which are not: the reliance of promisees, the needs of those wronged, and the adult's present values. The separation explains why some childhood obligations survive into adulthood with undiminished force and most do not.
\end{abstract}

\section{The principle and its scope}

PID-07-a, Principle~P1, states that the strength of an obligation grounded in the relation between a person at one time and herself at another is an increasing function of their psychological connectedness, with a small residual weight for identity. With Lemma~3.2, that connectedness between a mature adult and her self of middle childhood is weak, it yields Theorem~3.3: the obligations so grounded are weak.

The principle concerns one ground of obligation. Many obligations that arise from childhood acts have several grounds, of which the relation between the adult and the child is only one.

\begin{definition}\label{d:comp}
The \emph{agential component} of an obligation arising from an earlier act is the part of its force that depends on the obligated person's being the one who performed the act. Its \emph{independent components} are the parts that would bind her equally if someone else had performed the act and she stood in her present relations to its effects.
\end{definition}

\begin{proposition}\label{p:split}
The agential component of an obligation arising from a childhood act is weak; its independent components are not thereby weakened.
\end{proposition}

\begin{proof}
The agential component is grounded in the relation between the adult and her childhood self, so it falls under Principle~P1 and is weak by Theorem~3.3 of PID-07-a. The independent components are grounded in present relations, between the adult and other persons or the effects of the act, and the principle does not apply to them.
\end{proof}

\section{Three kinds of case}

\paragraph{Promises.} A girl of nine promises her younger brother that she will always look after his rabbit, or promises her dying grandfather that she will become a doctor. Two factors weaken such promises apart from identity: children's capacity to bind themselves is limited, and the content of the promise was understood by a mind very different from the adult's. Proposition~\ref{p:split} adds a third: the adult's link to the act of promising is weak. But promissory obligation has a further ground. On ETH-16, promises bind because they create reasonable reliance, and the obligation persists even when the promisee forgets; compare the role of expectation in the principle of fidelity of \citet[ch.~7]{Scanlon1998}. Where the promisee has reasonably relied on a childhood promise, as when the brother arranged his life around his sister's undertaking, the reliance is a present relation between the adult and the promisee, and its claim is not weakened by the weakness of her connection with the child. Most childhood promises create little reasonable reliance, since others know that children's promises are not to be relied on as an adult's are. The promise to the dying grandfather created none that survives him. So most childhood promises lapse or nearly lapse, and the exceptions are explained.

\paragraph{Wrongs.} A boy of ten bullies a classmate, who suffers lasting harm. Is the man of forty responsible? His responsibility in the sense of blameworthiness belongs to the agential component, and is weak: the attitudes and dispositions that issued in the bullying are, by hypothesis, largely gone, and blame directed at him for them is directed at someone only weakly connected with their source \citep[ch.~15]{Parfit1984}. Two independent components remain. The classmate's continuing harm gives reasons to anyone well placed to remedy it, and the adult is especially well placed, since an acknowledgment from him can do what one from no one else can. And there is a residual weight for identity: it was he, and not someone else, who did it, which gives him a standing to apologize and a reason to do so that a stranger lacks. The adult's obligation of repair is therefore real, and in serious cases substantial, though his culpability is slight.

\paragraph{The child's hopes.} The child hoped to become an astronaut, to live by the sea, never to work in an office. Does the adult wrong her childhood self by abandoning these hopes? The agential component is again weak. But hopes the adult has taken up as her own values connect her with her childhood self, since the persistence of values is a direct connection (PID-07-a, Definition~2.1), and to that extent the connection, and the reason, is strong. A childhood hope the adult still shares has a claim on her because she shares it. One she has outgrown has little. The child's hopes may also have a claim of another kind: the claim, which parents and guardians should respect, that the child's future be left open so that she can later choose among such hopes \citep{Feinberg1980}. That claim binds those who act for the child, not the adult she becomes.

\begin{objection}
The view lets adults escape what they did as children too easily. A man who as a teenager caused a death cannot shrug it off as the act of a weakly connected self.
\end{objection}

\begin{reply}
Theorem~3.3 of PID-07-a concerns the childhood self, the stages before adolescence. Connectedness with a teenage self is stronger, and the agential component correspondingly stronger. In any case the duty of repair is an independent component, and it is not reduced. What is reduced is the warrant for continuing to treat the man as the culpable agent of an act whose motivational sources are no longer his, which is a conclusion that the practice of treating childhood offenses differently already reflects.
\end{reply}

\section{Conclusion}

The weakness of the obligations grounded in the relation between an adult and her childhood self is the weakness of one component. Reliance, harm that continues and values that persist ground independent components, which explains why a sister may still owe her brother what she promised as a girl, why a man must still try to repair what he did as a boy, and why a child's hope binds the adult only if it is still hers.

\begin{thebibliography}{99}\small
\bibitem[Feinberg(1980)]{Feinberg1980} Feinberg, J. (1980). The child's right to an open future. In W. Aiken and H. LaFollette (eds.), \emph{Whose Child? Children's Rights, Parental Authority, and State Power}. Rowman and Littlefield.
\bibitem[Parfit(1984)]{Parfit1984} Parfit, D. (1984). \emph{Reasons and Persons}. Oxford University Press.
\bibitem[Scanlon(1998)]{Scanlon1998} Scanlon, T.~M. (1998). \emph{What We Owe to Each Other}. Harvard University Press.
\end{thebibliography}

\end{document}
